\section{Canonical ideals: reconstruction on the entire boundary}
\label{sec:canonical}
\textit{Support: the entire qualified G331 only. None of the current408 or historical notes is a correctness premise.}

Let $(R,\mathfrak m)$ be excellent local. Semivaluations kill nilpotents, so throughout this section replace $R$ by $R_{\rm red}$. If this reduction has dimension zero its centred boundary is empty and all envelopes mean its unique empty function. Otherwise fix a proper $\mathfrak m$-primary ideal $\mathfrak a$ and the compact, nonempty section $X=V_{\mathfrak a}(R)$. Every prime support strictly below $\mathfrak m$ is retained. G331 supplies the independent geometric core: actual negative effective finite real Cartier presentations on directed actual primary blowups, with degree nonnegative on every contracted proper integral curve. Its uniform class $C$ and actual decreasing-net class $D$ equal the corresponding ideal classes. We use this qualified comparison, its rational-profile density, finite maxima, full-fibre positive descent, and its exact published smooth specialization. No envelope theorem is part of that grant.

For a finite continuous signed obstacle $f:X\to\mathbb R$ define
\[
 A_0=R,\qquad A_m=\{g\in R:\ v(g)+mf(v)\ge0\text{ for all }v\in X\}\quad(m\ge1),
 \qquad h_m=-v(A_m)/m.
\]
This definition uses no fraction field of $R$ and is meaningful also when $v(g)=+\infty$. Set $c=\max(0,-\min_X f)$. Any positive integer $N\ge mc$ gives $\mathfrak a^N\subset A_m$. The ultrametric inequality proves addition closure and nonnegative values on $R$ prove multiplication closure. If $f$ is negative at one point, no unit lies in any $A_m$, so every $A_m$ is proper primary. If $f\ge0$, every $A_m=R$. Thus $h_m$ is always a finite continuous unit-or-primary profile, is below $f$, and has actual geometric positivity by G331. In the proper case its presentation is $-E_{A_m}/m$ on the actual blowup of $A_m$; the negative exceptional bundle is relatively ample.

Multiplying the defining inequalities gives $A_mA_n\subset A_{m+n}$. Each $A_m$ is integrally closed as an ideal. Indeed, if $g$ satisfies a monic degree-$q$ equation with coefficient of index $i$ in $A_m^i$, then at a semivaluation where $v(g)<v(A_m)$ the leading term has uniquely least finite value. This contradicts the equation. All ideal values here are finite because $A_m$ contains a primary power. Infinite $v(g)$ imposes no difficulty.

Define only a raw function, without a positivity assertion,
\[
 p(f)=\sup\{u\in D:u\le f\}.
\]
We prove
\[
 p(f)=\sup\{u\in C:u\le f\}=\sup_{m\ge1}h_m.
 \tag{1}
\]
A rational unit-or-primary competitor $-v(I)/r\le f$ has $I\subset A_r$, and hence is below $h_r$. A $C$ competitor can be approximated uniformly by rational profiles and then shifted down by a rational negative $\mathfrak a$-constant. Such a shift is again a rational primary profile, by multiplication of ideals after clearing denominators. Sending the approximation error to zero proves the $C$ assertion. If $u=\inf_i g_i\le f$ in an actual decreasing $C$-net, the open sets $\{g_i<f+\epsilon\}$ cover $X$. Directedness and compactness select one index valid everywhere. The function $g_i-\epsilon\in C$ is a competitor and $u\le\sup_m h_m+\epsilon$. Let $\epsilon\downarrow0$. This step uses continuity of $f$, not continuity of $u$; bottom and partial poles are permitted. The reverse inequalities follow because each $h_m$ is itself a $D$ competitor.

Put $M_1=\|v(A_1)\|_\infty$ and $H_N=\max_{1\le r\le N}h_r\in C$. Write $m=kr+s$, $0\le s<r$. The inclusion $A_r^kA_1^s\subset A_m$, together with $h_r\le0$, gives
\[
 h_m\ge \frac{kr}{m}h_r-\frac{s}{m}M_1
 \ge h_r-\frac{rM_1}{m},\qquad
 H_N-\frac{NM_1}{m}\le h_m\le p(f).
 \tag{2}
\]
These estimates hold on every support, even when generator values are infinite. In particular the full degree sequence converges pointwise to $p(f)$. The function $p(f)$ is finite, nonpositive and lower semicontinuous, bounded below by $h_1$.

The following exact criterion is not a solution of general continuity:
\[
 p(f)\text{ continuous}\quad\Longleftrightarrow\quad H_N\to p(f)\text{ uniformly}
 \quad\Longleftrightarrow\quad h_m\to p(f)\text{ uniformly}. \tag{3}
\]
For the forward implication the increasing strict error-test open sets for $H_N$ cover compact $X$, so one $N$ works uniformly. Equation (2) then gives uniform convergence of the full sequence. Conversely a uniform limit lies in the closed geometric class $C$, so is continuous and is the greatest $D$ minorant by (1). We do not declare an upper semicontinuous regularization positive.

The raw envelope is monotone and satisfies
\[
 \|p(f)-p(g)\|_\infty\le\|f-g\|_\infty. \tag{4}
\]
If $f\le g+\epsilon$, any $D$ minorant $u$ of $f$ has $u-\epsilon\in D$ below $g$. The shift belongs to $D$ by adding a fixed negative primary profile to its defining decreasing net. Supremizing proves the one-sided inequality; reversing $f,g$ proves (4). Consequently uniform approximation by obstacles whose envelopes are already in $C$ proves a $C$ envelope for the limit. This is an implication, not a general density assertion.

For completeness we independently reconstruct the finite-generation reduction. Suppose $\mathcal A=\bigoplus_{m\ge0}A_mt^m$ is generated by homogeneous $g_it^{d_i}$, $d_i>0$. Let $L$ be a common multiple of the degrees, $D_0=\max d_i$, and $J=(g_i^{L/d_i})_i$. On each support some generator has finite value: otherwise the degree-one ideal $A_1$ would have infinite value. Every degree-$m$ expression has value at least $m\min_i v(g_i)/d_i$. Moreover $J\subset A_L$, so
\[
 h_m\le -v(J)/L\le h_L\le p(f).
\]
The middle profile is therefore $p(f)=h_L$, and it is also $\max_{r\le D_0}h_r$ by using a generator attaining the finite minimum. The radical of the generator coefficient ideal contains $A_1$, so $J$ is primary in the proper case. Equation (2) gives
\[
 0\le p(f)-h_m\le D_0M_1/m. \tag{5}
\]
The unit case is immediate.

Conversely, if $p(f)=-v(J)/q$ for a proper primary $J$ and positive integer $q$, then the definition and (1) imply
\[
 A_m=\{g:v(g)\ge (m/q)v(J)\text{ on all }X\}
     =\{g:g^q\in\overline{J^m}\}. \tag{6}
\]
For the second equality apply the centred rank-one DVR integral-closure criterion to each minimal-component quotient and then the all-minimal-component criterion (Swanson--Huneke 6.8.3--4 and 1.1.5). Every such quotient has positive dimension: a zero-dimensional minimal component of a local reduced ring would make its maximal ideal minimal and the whole ring zero-dimensional. The images of $J$ and $\mathfrak a$ are primary, and each centred DVR normalizes into $X$. Monic ideal-integrality equations give the reverse inequalities on all other supports as well, by the same uniquely least finite term argument. These are applications of the complete immediate published units, read afresh; their deeper valuation-domination foundations are disclosed in the ledger.

Set $B=R[Jt^q]$. Primary $J$ contains a regular element by finite prime avoidance. Thus $B$ is reduced excellent and its total quotient ring is $\prod_j K_j(t^q)$, with $K_j=\operatorname{Frac}(R/\mathfrak p_j)$. Its finite extension $\prod_j K_j(t)$ may be inseparable. The standard excellent/Nagata finite-normalization property makes the integral closure of $B$ in that extension finite over $B$. An equation for $g^q$ over $J^m$, multiplied by powers of $t^{mq}$, makes $g^qt^{mq}$ integral over $B$, and then $gt^m$ integral by its $q$th-power equation. Hence $B\subset\mathcal A$ and $\mathcal A$ is a $B$-submodule of that finite normalization. Noetherianity makes it finite over $B$, hence finitely generated over $R$. The same statement for the zero envelope is $\mathcal A=R[t]$. A finite maximum of canonical profiles is a rational unit-or-primary profile, by ideal sums after clearing denominators. We have proved equivalence of finite generation, rational primary attainment, and finite canonical-max attainment, without assuming any of them generally.

For a signed integral exceptional Cartier obstacle $f_D$ on an actual primary blowup of a normal $R$, let $Y$ be its finite normalization, with the pulled-back divisor and its canonical rational trivialization. There is an exact model computation
\[
 A_m(f_D)=H^0(Y,\mathcal O_Y(mD))\cap R. \tag{6a}
\]
Testing the all-boundary inequalities at the finitely many height-one exceptional orders gives the section inequalities there; at other height-one primes $D$ is zero and $g\in R$ is already regular. Normal height-one intersection proves that $g$ is a section. Conversely a section pulls back to the valuation ring of $\operatorname{Frac}(R/\mathfrak p)$ at every semivaluation support. The generic point is outside the closed exceptional fibre, so the invertible bundle has its canonical nonzero trivialization there. If its local generator is $b$, then the section has $v(g)-v(b)\ge0$, and $v(b)=-mf_D(v)$. This yields the canonical inequality on every support, with infinite $v(g)$ automatic. Properness gives the valued centre; no total quotient field is mapped through a nonminimal support. For a rational divisor use denominator multiples. The finite generation of this particular section-intersection algebra is a hypothesis to establish in any application, not a theorem for arbitrary signed model obstacles.

For another primary gauge $\mathfrak b$, put $d(v)=v(\mathfrak b)>0$ and $T(v)=v/d(v)$. If $f_{\mathfrak b}(Tv)=f_{\mathfrak a}(v)/d(v)$, then
\[
 (Tv)(g)+mf_{\mathfrak b}(Tv)=(v(g)+mf_{\mathfrak a}(v))/d(v).
\]
The canonical ideals and graded algebra are literally unchanged. Profiles and raw envelopes divide by $d$; substitution gives the three-gauge cocycle. Uniform errors divide by the positive minimum of $d$. For the finite normalization $S$ of $R$, surjective restriction gives $A_m^R=A_m^S\cap R$. This includes every conductor-infinite quotient fibre; it does not assert that contraction preserves finite generation. If the source geometric envelope exists and is constant on every normalization fibre, G331's positive descent supplies the downstairs envelope: it is below $f$ and dominates every pulled-back downstairs competitor. Compatibility of the obstacle alone is not sufficient.

If a decreasing directed family of continuous obstacles has envelopes $H_i\in C$, then $H=\inf_iH_i\in D$ is the envelope of their infimum: every minorant of the infimum is below every $H_i$, and $H$ is below every obstacle. This uses an actual decreasing net and no flattening. For excellent reduced curves the finite normalization centres are DVRs. A regular conductor element has finite order on each, so all boundary fibres are singleton. The independent geometric DVR core has every finite nonpositive real slope, with no contracted curves. Positive descent yields the entire nonpositive finite orthant; decreasing truncations yield the extended orthant. Consequently $P(f)=\min(f,0)$ geometrically, at every branch. The literal classical/CLD multibranch radial class still excludes mixed-zero profiles and is not identified with this boundary class.

On a DVR, $f=-\alpha$ with $\alpha>0$ irrational gives $A_m=(x^{\lceil m\alpha\rceil})$ and the continuous real Cartier envelope $-\alpha$. Finite homogeneous generators would force a rational minimum order/degree ratio, contradicting the limiting slope $\alpha$. Thus finite generation is not necessary for continuity. This diagnoses the limitation of (5), rather than an envelope failure.
